\section{Related Work}

\subsection{Algorithmic Problem Solving Education}

Algorithmic problem solving has been considered a fundamental skill in
computer science education and software engineering recruitment.

Traditional approaches emphasize the understanding of data structures,
algorithm design techniques, and computational complexity analysis.
Platforms such as LeetCode have become widely adopted tools for students
to practice programming problems and prepare for technical interviews.

Previous educational research has investigated how repeated programming
practice improves algorithmic thinking, pattern recognition, and problem
solving ability.

However, traditional algorithm training requires significant time investment,
as students often solve hundreds of problems to develop sufficient experience.


\subsection{Large Language Models for Code Generation}

Recent advances in Large Language Models (LLMs) have significantly changed
software development workflows.

Models such as GPT-based systems and other code-capable LLMs have demonstrated
the ability to generate code, explain algorithms, and assist developers in
debugging tasks.

Previous studies have evaluated LLM performance on programming benchmarks and
found that these models can solve many software engineering tasks while still
facing challenges in correctness, reasoning, and complex algorithm selection.

These limitations suggest that effective interaction strategies between humans
and LLMs remain an important research direction.


\subsection{Prompt Engineering for Software Development}

Prompt engineering has emerged as a method to improve LLM performance by
providing structured instructions and reasoning guidance.

Research has shown that carefully designed prompts can improve model outputs
in tasks involving reasoning, planning, and code generation.

In software engineering contexts, prompts that include requirements analysis,
algorithm planning, and verification steps may help LLMs produce higher-quality
solutions.

However, limited research has investigated whether prompt engineering can
change the role of traditional algorithm training, especially in the context
of software engineering interview preparation.


\subsection{Research Gap}

Existing studies have primarily focused on evaluating whether LLMs can generate
correct code or improve developer productivity.

However, fewer studies have examined the broader question:

\begin{quote}
Will algorithmic training practices for future software engineers shift from
traditional problem solving toward AI-assisted prompt engineering workflows?
\end{quote}


This research addresses this gap by comparing different prompting strategies
for algorithmic problem solving and investigating how AI-assisted programming
may influence future software engineering skill requirements.
